\section{Experiments}
\label{sec:experiments}

\Cref{sec:analysis} bounds the forward error of a single reduction and
decomposes its effect on the loss at the output softmax. Neither result
predicts the perplexity ordering of SR and RN once rounding errors propagate
through the network: the bounds are worst-case and per reduction, and the loss
decomposition is evaluated one site at a time. This section measures that
ordering on DistilGPT-2 at matched significand precision, with the scope of
reduced precision ranging from the whole network to one projection in one block
(\cref{tab:sweeps}). The sweeps address three questions: which rounding rule
gives lower perplexity at each site, how per-site effects combine across sites
and across depth, and whether the per-site preferences yield a mixed-rounding
assignment that improves on RN at identical per-site precisions.

\subsection{Experimental setup}
\label{sec:exp-setup}

\paragraph{Model and evaluation benchmark.}
We evaluate the public \code{distilbert/distilgpt2} checkpoint
\citep{distilgpt2modelcard,wolf2020transformers}. The architecture retains the
layer widths of GPT-2 small ($d_{\rm model} = 768$, $d_{\rm ff} = 3072$, $H =
  12$ heads, $d_{\rm head} = 64$, vocabulary $V = 50257$, as illustrated in
\cref{fig:architecture}) and distills the depth to six blocks. It therefore has
the reduction lengths that enter the bounds of \cref{sec:analysis-mlp}:
$n=768$ for \code{attn\_c\_attn}, \code{attn\_c\_proj}, and \code{mlp\_c\_fc},
and $n=3072$ for \code{mlp\_c\_proj}. Six blocks keep per-operation software
emulation tractable. All evaluations use the pretrained weights without
fine-tuning, calibration, or quantization-aware adjustment, so that
configurations differ only in the arithmetic under test.

Evaluation uses the WikiText-2 test split \citep{merity2017pointer}. The harness
selects the leading $1\%$ of lines, joins them with double newlines, tokenizes
the resulting text, and evaluates its first $1024$ tokens. These are divided
into four non-overlapping contexts of $256$ tokens, with context reset between
them. Each context contributes $255$ next-token targets, giving $1020$ scored
positions. Perplexity is the exponential of the mean cross-entropy over these
positions; the PRISM RN reference at \code{binary32} significand precision
($t=24$) is $\mathrm{PPL}_0 = 62.51$. When effects are compared or summed, we
also report the excess loss $\ell = \ln(\mathrm{PPL}/\mathrm{PPL}_0)$ in nats,
which equals the increase in mean cross-entropy over the $1020$ positions and is
the quantity decomposed in \cref{eq:head-full-comparison}.

\paragraph{Instrumented arithmetic and scoping.}
\label{sec:exp-scoping}
Inference is executed within Fuzzy PyTorch \citep{gonzalezpepe2026fuzzy}. During
the framework's compilation, the Verificarlo compiler
pass~\citep{denis2016verificarlo} intercepts native floating-point operations
and redirects them to PRISM's vectorized C++ backend (\cref{sec:prism}).
Instrumented operations within the target scope are rounded to a virtual
significand precision $t \le 24$ using either SR or untied RN, in which half-ULP
ties round away from zero rather than to even (\cref{sec:prism}). All SR--RN
comparisons in this section use untied RN and not IEEE round-to-nearest with
ties-to-even (RNE) at reduced precision. The two differ only on exact ties. RN
introduces no randomness, and repeated RN runs are bitwise reproducible.

The harness attaches PyTorch forward pre-hooks that set precision $t$ and the
rounding mode immediately before the target module executes, and forward
post-hooks that restore $t = 24$ under RN on exit. Inference runs in PyTorch
eager mode without cross-module operator fusion. Every operation outside the
target scope, including the final loss and perplexity computation, runs at
$t=24$ under RN, so an SR run and its paired RN control differ only in the
rounding rule applied within the target scope.

The component and sublayer sweeps reduce precision at the target in all six
blocks at once; the depth sweeps reduce it in one block or in a prefix of
blocks (\cref{tab:sweeps}). The \code{attention} and \code{mlp} scopes cover
every operation inside the corresponding module of each block: the query--key--value
and output projections, the score and value products, and the softmax for
\code{attention}; \code{mlp\_c\_fc}, the GELU, and \code{mlp\_c\_proj} for
\code{mlp}. The \code{attn\_qkav} scope covers only the products $QK^\top$
(reduction length $d_{\rm head}=64$) and $AV$ (at most $T=256$), with the
softmax at $t=24$. The \code{lm\_head} scope is the projection onto the
vocabulary logits.

\begin{table}[t]
  \centering
  \caption{Sweep levels. All levels compare SR and untied RN at matched
    significand precision $t$; operations outside the target scope run at
    $t=24$ under RN.}
  \label{tab:sweeps}
  \small
  \begin{tabular}{@{}l >{\raggedright\arraybackslash}p{4.2cm} c >{\raggedright\arraybackslash}p{4.8cm}@{}}
    \toprule
    Level & Target scope                                                                  & $t$ swept & Question \\
    \midrule
    Global
          & Whole model
          & $4\text{--}14$
          & SR vs. RN comparison under uniform precision                                                         \\
    Component
          & \code{attention}, \code{mlp}, \code{lm\_head}, \code{attn\_qkav} (all blocks)
          & $4\text{--}14$
          & Site-dependent sensitivity across main modules                                                       \\
    Sublayer
          & Four projections \code{attn\_*}, \code{mlp\_*} (all blocks)
          & $4\text{--}14$
          & Impact of reduction length ($n=768$ vs.\ $3072$)                                                     \\
    Depth
          & \code{mlp\_c\_proj} (single blocks; prefixes $0\text{--}k$)
          & $4,5,6$
          & Dependence on position; additivity of per-block effects                                              \\
    Mixed recipe
          & SR at output projections, RN elsewhere
          & $6,7,8$
          & Mixed rounding versus all-RN at identical per-site precisions                                        \\
    \bottomrule
  \end{tabular}
\end{table}

\paragraph{Statistical protocol.}
The global, component, sublayer, and depth sweeps use five independent seeds per
SR configuration, and the mixed-rounding comparison in
\cref{sec:exp-composition} uses five seeds per mixed configuration. Results are
reported as sample mean $\pm$ standard deviation; for SR, excess loss is
averaged over seeds. RN is deterministic at a fixed execution configuration and
is run once. We define the perplexity gap as $\Delta = \overline{\SR} - \RN$, so
negative values favor SR. We regard the gap as statistically distinguishable from
zero when the RN value falls outside the two-sided $95\%$ Student-$t$ confidence
interval for the SR mean ($df=4$). The evaluated text is fixed, so every
interval is conditional on the $1020$ scored positions and does not describe
variation across texts.

\subsection{Results}
\label{sec:exp-results}

The experiments yield four findings:
\begin{enumerate}[label=(\roman*),leftmargin=2em]
  \item \textbf{RN is better under uniform rounding.}
        Reducing only the language-model head reproduces most of the
        SR--RN gap, suggesting that the head drives the global preference
        for RN (\cref{fig:global,fig:component}).

  \item \textbf{The preferred rounding rule depends on the site.}
        SR gives lower perplexity in the MLP, whereas RN is preferred at
        the head and attention input projection. The loss decomposition
        attributes the MLP advantage to reduced drift and the head
        disadvantage to SR fluctuations
        (\cref{fig:component,fig:sublayer}).

  \item \textbf{RN loss compounds more strongly across blocks.}
        For \code{mlp\_c\_proj} at $t=5$--$6$, RN excess loss exceeds
        the sum of single-block losses more than SR excess loss does.
        At $t=6$, SR is approximately additive
        (\cref{fig:cumulative}).

  \item \textbf{Mixed-rounding improves on full RN.}
        Using SR at the attention and MLP output projections and RN
        elsewhere lowers perplexity, even with one fewer significand
        bit at the MLP down-projection
        (\cref{tab:accuracy-matched-precision}).
\end{enumerate}

\paragraph{Global uniform-precision sweep.}
Under uniform precision, RN consistently outperforms SR across the entire
low-to-moderate precision range, with statistically significant differences at
every $t \in [4,11]$ (\cref{fig:global}). At $t=8$, perplexity is $97.68$ under
RN and $139.27\pm3.58$ under SR. Reducing only the head reproduces this gap: the
SR--RN difference in excess loss is $2.03$ nats globally and $2.27$ nats with
only \code{lm\_head} reduced at $t=6$, and $0.35$ against $0.37$ nats at $t=8$.
The RN advantage is therefore consistent with the head's SR fluctuation penalty
(\cref{sec:analysis-empirical-drift}) dominating the whole-network result. This
general performance agrees with the forward-pass preference for RN reported by
\citet{chmiel2023accurate}, establishing RN as the appropriate default for
global precision reduction.

\begin{figure}[t]
  \centering
  \safeincludegraphics{figures/global_sweep.pdf}
  \caption{Global sweep at uniform precision across all model operations, showing mean perplexity and SR standard deviation on a logarithmic axis. RN's advantage is statistically significant across $t\in[4,11]$.}
  \label{fig:global}
\end{figure}

\begin{figure}[t]
  \centering
  \safeincludegraphics{figures/component_sweep.pdf}
  \caption{Per-component sweep across precisions $t \in [4, 14]$. Bands show $\pm 1$ sample SD over five SR seeds; dotted line indicates the full-precision reference ($62.51$).}
  \label{fig:component}
\end{figure}

\paragraph{Site performance: MLP versus language-model head.}
In the MLP, SR gives lower perplexity than RN at every $t\le8$, with
statistically significant differences (\cref{fig:component}). At $t=6$,
perplexity is $71.97\pm1.16$ under SR and $137.93$ under RN, a degradation of
$15\%$ against $121\%$ over the reference. At $t=8$, SR is within $0.3\%$ of the
reference ($62.69$) and RN within $2.0\%$ ($63.75$). At $t\le5$ both rules more
than double the reference perplexity (at $t=5$, $175.65$ under SR against
$2029.83$ under RN); the ordering there is statistically significant but
concerns configurations that neither rule makes usable.

At the language-model head the ordering reverses: RN gives lower perplexity at
every $t\in[4,12]$, with statistically significant differences. At $t=8$,
perplexity is $135.56\pm5.22$ under SR and $93.60$ under RN; at $t=6$,
$5054\pm464$ against $519.18$. \Cref{sec:analysis-empirical-drift} decomposes
the head's excess loss at these precisions. Over $t=6,\ldots,9$ the fluctuation
term $V_{\SR}$ accounts for $100\%$ of SR's excess loss, whereas RN,
which has no fluctuation, pays only for its non-uniform drift ($2.12$ against
$4.40$ nats at $t=6$). SR noise injected at the head changes the differences
between logits, which the cross-entropy penalizes; SR noise injected in the MLP
reaches the logits predominantly as a uniform shift, to which the softmax is
invariant.

The \code{attention} module favors SR at $t\le6$ but do not favor any rounding
mode at higher precisions; at $t\ge7$ the two rules differ by less than $2\%$ of
the reference (at $t=7$, $64.60\pm0.49$ under SR against $63.42$ under RN).
Restricted to the score and value products (\code{attn\_qkav}), reduced
precision raises perplexity by at most $14\%$ at $t=4$ ($71.26\pm1.26$ under SR,
$65.90$ under RN) and by at most $2.9\%$ at $t\ge5$. These products have the
shortest reductions in the model ($d_{\rm head}=64$ and at most $T=256$ terms),
which may explain their small sensitivity to either rule.

The \code{attention} module does not decompose into its parts. At $t=4$, the
excess losses of \code{attn\_c\_attn}, \code{attn\_c\_proj}, and
\code{attn\_qkav}, each reduced separately, sum to $1.20$ nats under RN and
$1.66$ nats under SR, which favors RN; reducing the whole module gives $2.56$
and $2.29$ nats, which favors SR. At $t=6$ the sums are equal under both rules
($0.094$ nats), whereas the whole module gives $0.182$ under RN and $0.120$
under SR.

\begin{figure}[t]
  \centering
  \safeincludegraphics{figures/sublayer_sweep.pdf}
  \caption{Per-sublayer sweep across the four primary projection matrices, with
    reduction lengths $n$ indicated and perplexity on a shared logarithmic scale.
    The feedforward down-projection (\code{mlp\_c\_proj}, $n=3072$) exhibits the
    largest low-precision SR advantage ($\Delta = -1473.15$ at $t=4$).}
  \label{fig:sublayer}
\end{figure}

\paragraph{Reduction length and sublayer sensitivity.}
\Cref{fig:sublayer} separates the four projections, each reduced in all six
blocks. The MLP down-projection (\code{mlp\_c\_proj}, $n=3072$) has the largest
SR advantage, statistically significant at every $t\le7$: at $t=6$, perplexity
is $66.24\pm0.63$ under SR and $79.64$ under RN, a degradation of $6.0\%$
against $27\%$. Among the three projections with $n=768$, the preferred rule
depends on the site. The attention input projection (\code{attn\_c\_attn})
favors RN wherever the difference is statistically significant ($t=4,5,7$; at
$t=4$, $251.58\pm16.05$ under SR against $154.08$ under RN). The attention
output projection (\code{attn\_c\_proj}) favors SR at $t=4$--$6$, and the MLP
input projection (\code{mlp\_c\_fc}) favors SR at $t=5$--$6$; at $t\ge7$ neither
difference is statistically significant.

The largest SR advantage occurring at the longest reduction is consistent with
\cref{sec:analysis-mlp}: RN rounding errors can align across the $n$ terms of a
reduction, whereas SR errors have zero conditional mean, so RN drift can grow
faster with $n$ than SR fluctuation. The sweep does not isolate $n$, however.
Only one site has $n=3072$, and it also differs from the others in its input,
the GELU output, which is bounded below by approximately $-0.17$. The opposite
preferences of three sites with identical $n=768$ show that factors other than
$n$, such as the input distribution, the weights, and the downstream
operations, also determine the ordering.


\paragraph{SR's advantage across depth.}
The cumulative sweep in \cref{fig:cumulative} reduces \code{mlp\_c\_proj} in
progressively longer block prefixes, from block~$0$ to all six blocks. For every
tested precision, SR yields statistically significantly lower perplexity than RN
at every prefix, except $0$ and $0$--$1$ at $t=6$, where the differences are not
statistically significant.

The advantage of SR becomes more pronounced as more blocks use reduced
precision. At $t=4$, perplexities are initially close ($83.24$ under SR and
$85.13$ under RN for block~$0$), but extending the reduction to all six blocks
increases perplexity by a factor of $5.5$ under SR and $22.7$ under RN. At
$t=6$, the difference becomes statistically significant from prefix $0$--$2$
onward, reaching $66.24\pm0.63$ under SR versus $79.64$ under RN with all six
blocks reduced, as in \cref{fig:sublayer}. These results show that SR limits the
perplexity degradation from reducing the MLP down-projection across multiple
blocks. The growing SR--RN gap is consistent with RN introducing systematic drift
and SR limiting its accumulation across blocks.

This compounding effect illustrates a key structural difference between the
rounding rules. The deterministic drift introduced by RN in one block acts as a
correlated error that systematically shifts activations in subsequent blocks. In
contrast, the noise introduced by SR accumulates without compounding directional
drift across depth. While SR's unbiasedness does not survive non-linear
operations, empirically, the downstream propagated noise remains far less biased
for SR than for RN, validating its theoretical benefits across network depth.

\begin{figure}[t]
  \centering
  \safeincludegraphics{figures/cumulative_propagation.pdf}
  \caption{Cumulative propagation across block prefixes $0$ to $0\text{--}k$ for \code{mlp\_c\_proj} at $t \in \{4, 5, 6\}$. Bands show $\pm 1$ sample SD over five SR seeds.}
  \label{fig:cumulative}
\end{figure}


\subsection{Mixed rounding beats full RN at matched precision}
\label{sec:exp-composition}

The sweeps above reduce one site or one module at a time, and per-site effects
do not add (\cref{sec:exp-results}). We therefore test whether the per-site
preferences hold when five sites are reduced together. Each mixed configuration
assigns SR or RN per site and is compared with an all-RN configuration at
identical per-site precisions (\cref{tab:accuracy-matched-precision}).

The assignment follows the sublayer sweeps. SR is used at the two output
projections, where it is preferred in isolation: \code{attn\_c\_proj} at $t=6$
($62.10\pm0.42$ under SR against $62.91$ under RN) and \code{mlp\_c\_proj} at
$k\in\{6,7,8\}$ (statistically significant at $k=6,7$; not statistically
significant at $k=8$). RN is used at the two input projections at $t=8$, where
neither sublayer difference is statistically significant and
\code{attn\_c\_attn} favors RN at every lower precision with a statistically
significant difference. The head runs under RN at $t=11$, where head-only RN is
within $0.8\%$ of the reference ($63.00$). All remaining operations run at
$t=24$ under RN. The all-RN control replaces SR by RN at the two output
projections and changes nothing else.

\begin{table}[t]
  \centering
  \caption{Mixed rounding versus RN at identical per-site significand
    precisions on WikiText-2. Lower perplexity is better.}
  \label{tab:accuracy-matched-precision}
  \footnotesize
  \setlength{\tabcolsep}{5pt}
  \renewcommand{\arraystretch}{1.15}
  \begin{tabular}{@{}l c cc@{}}
    \toprule
    \multicolumn{3}{@{}l}{\textbf{(a) Precision and rounding configurations}}                                \\
    Site                                    & $t$ bits         & Mixed                             \\
    \midrule
    Attention input (\code{attn\_c\_attn})  & $8$              & RN                                \\
    Attention output (\code{attn\_c\_proj}) & $6$              & \textbf{SR}                       \\
    MLP input (\code{mlp\_c\_fc})           & $8$              & RN                                \\
    MLP output (\code{mlp\_c\_proj})        & $k \in \{6,7,8\} $          & \textbf{SR}            \\
    Language-model head (\code{lm\_head})   & $11$             & RN                                \\
    All other operations                    & $24$             & RN                                \\
    \bottomrule
  \end{tabular}\hfill
  \begin{minipage}[c]{0.42\linewidth}
    \footnotesize
    Configurations apply across all six blocks; $\mu\pm\sigma$ over five seeds.
    Negative $\Delta$ favors mixed rounding. CIs are $95\%$ Student-$t$ ($df=4$).
    Relative baseline column shows percentage increase over the full-precision reference ($62.51$) for Mixed (Full RN).
  \end{minipage}

  \medskip
  \begin{tabular}{@{}c c c c c c@{}}
    \toprule
    \multicolumn{6}{@{}l}{\textbf{(b) Perplexity gaps}}                                                                                                          \\
    \addlinespace[2pt]
    \shortstack{$k$}
        & \shortstack[r]{Mixed ($\mu\pm\sigma$) }
        & \shortstack[r]{Full RN}
        & \shortstack[r]{Gap $\Delta$ ($\%$)}
        & \shortstack{$95\%$ CI $\Delta$}
        & \shortstack[r]{Rel. baseline $\Delta$ (Full RN)}                                                                                                       \\
    \midrule
    $8$ & $64.63\pm0.68$                                   & $66.66$ & \SRbetter{$-2.03$ ($-3.04\%$)}   & $[-2.87,-1.18]$   & \RNbetter{$+3.38\%$ $(+6.64\%)$}   \\
    $7$ & $64.89\pm0.98$                                   & $68.89$ & \SRbetter{$-4.00$ ($-5.81\%$)}   & $[-5.22,-2.79]$   & \RNbetter{$+3.80\%$ $(+10.20\%)$}  \\
    $6$ & $69.01\pm2.07$                                   & $96.07$ & \SRbetter{$-27.06$ ($-28.17\%$)} & $[-29.63,-24.49]$ & \RNbetter{$+10.39\%$ $(+53.68\%)$} \\
    \bottomrule
  \end{tabular}
\end{table}

These results confirm that combining mixed rounding with mixed precision
consistently improves overall performance compared to using RN alone. At
identical per-site precisions, the mixed configuration gives lower perplexity
than all-RN from $3.04\%$ to $28.17\%$ for all three values of $k$. The gain
grows as the MLP down-projection loses precision, as in the sublayer sweep.
Across rows, the mixed configuration at $k=6$ ($69.01\pm2.07$) is not
distinguishable from all-RN at $k=7$ ($68.89$), and the mixed configuration at
$k=7$ ($64.89\pm0.98$) is below all-RN at $k=8$ ($66.66$), which lies outside
its $95\%$ interval. In this recipe, replacing RN by SR at the two output
projections compensates for removing one significand bit at the MLP
down-projection.

\subsection{Limitations}
\label{sec:exp-limitations}

PRISM emulates reduced significand precision while retaining the \code{binary32}
exponent range. These results concern significand and accumulator precision
under emulation; they do not model the exponent-range restrictions of narrow
hardware formats such as FP8 or NVFP4. Transfer to native hardware also depends
on its reduction order and arithmetic implementation. Non-target operations in
these sweeps use PRISM's untied RN at $t=24$, rather than native ties-to-even
(RNE); while exact tie events are rare in continuous floating-point activations, their frequency and impact have not been characterized.
Incorporating explicit exponent-range limitations remains future work.

The complete sweeps use one six-block model and one leading text prefix. The
measured site ordering and composition effects need replication across texts,
larger models, and broader tasks. The reduction bounds and loss decomposition
identify relevant quantities, but do not establish a universal rounding
assignment or quantitative portability of the observed gains.
